\section{Análisis de Sensibilidad de Parámetros}

Se realizó un análisis exhaustivo del impacto de los parámetros del algoritmo HO en las métricas de rendimiento para QC-DVRP. Los rangos evaluados, basados en Amiri et al. (2024), fueron:

\begin{itemize}
\item $\alpha \in [0.1, 0.9]$: Controla la agresividad en la fase de defensa
\item $\beta \in [0.2, 0.8]$: Modula la respuesta territorial
\item $\gamma \in [0.3, 1.0]$: Factor de evasión de amenazas
\end{itemize}

\subsection{Resultados Principales}

\begin{figure}[htbp]
\centering
\includegraphics[width=\textwidth]{parameter_sensitivity.pdf}
\caption{Efecto de los parámetros HO en métricas objetivo}
\label{fig:sensitivity}
\end{figure}

\begin{figure}[htbp]
\centering
\includegraphics[width=0.8\textwidth]{parameter_heatmap.pdf}
\caption{Mapa de calor: interacción $\alpha$-$\beta$ en entregas a tiempo}
\label{fig:heatmap}
\end{figure}

Los resultados indican que:
\begin{enumerate}
\item El parámetro $\alpha$ tiene el mayor impacto en la tasa de entregas a tiempo
\item Valores intermedios de $\beta$ (0.4-0.6) producen mejor balance
\item $\gamma$ afecta principalmente la capacidad de escape de óptimos locales
\end{enumerate}
